% Generado por sistema/backend ReporteTex.java -- no editar a mano.
% Datos: Base de datos: 148 importado (titulados) / Base de datos: 53 importado (empleadores)
\subsubsection*{C7. Financiamiento $\times$ nivel del programa cursado (opcional)}

\textit{Propósito.} Determinar si la fuente de financiamiento depende del nivel del programa cursado; complementa la historia de usuario de financiamiento. Respalda el criterio ARCU-SUR sin indicador ARCU-SUR directo. Se aplica sobre: con formación complementaria, con nivel de significancia $\alpha = 0.05$.

\textit{Paso 1. Contar con la muestra: tabla de contingencia} (frecuencia observada $FO_{ij}$, $N = 98$).

\begin{mitabla}{C7: frecuencia observada $FO_{ij}$}
\begin{tabular}{lcc|c}
\toprule
Financiamiento & Diplomado o Especialidad & Maestría o superior & Total fila \\
\midrule
Recursos propios & 66 & 24 & 90 \\
Beca total & 1 & 2 & 3 \\
Beca parcial & 2 & 3 & 5 \\
\midrule
Total columna & 69 & 29 & $N=98$ \\
\bottomrule
\end{tabular}
\end{mitabla}

\textit{Paso 2. Calcular la frecuencia esperada} de cada celda, $FE_{ij} = (\text{total fila } i)(\text{total columna } j)/N$. Por ejemplo, para la celda ``Recursos propios -- Diplomado o Especialidad'': $FE_{11} = 90 \times 69 / 98 = 63.3673$.

\begin{mitabla}{C7: frecuencia esperada $FE_{ij}$}
\begin{tabular}{lcc}
\toprule
Financiamiento & Diplomado o Especialidad & Maestría o superior \\
\midrule
Recursos propios & $90\cdot69/98 = 63.37$ & $90\cdot29/98 = 26.63$ \\
Beca total & $3\cdot69/98 = 2.11$ & $3\cdot29/98 = 0.89$ \\
Beca parcial & $5\cdot69/98 = 3.52$ & $5\cdot29/98 = 1.48$ \\
\bottomrule
\end{tabular}
\end{mitabla}

\textit{Paso 3. Calcular el estadístico chi-cuadrado,} sumando $(FO_{ij}-FE_{ij})^2/FE_{ij}$ de cada celda.

\begin{mitabla}{C7: cálculo de $X_0^2$ celda por celda}
\begin{tabular}{L{5.4cm}cccc}
\toprule
Celda & $FO_{ij}$ & $FE_{ij}$ & $(FO_{ij}-FE_{ij})^2$ & $(FO_{ij}-FE_{ij})^2/FE_{ij}$ \\
\midrule
Recursos propios -- Diplomado o Especialidad & 66 & 63.3673 & 6.9309 & 0.1094 \\
Recursos propios -- Maestría o superior & 24 & 26.6327 & 6.9309 & 0.2602 \\
Beca total -- Diplomado o Especialidad & 1 & 2.1122 & 1.2371 & 0.5857 \\
Beca total -- Maestría o superior & 2 & 0.8878 & 1.2371 & 1.3935 \\
Beca parcial -- Diplomado o Especialidad & 2 & 3.5204 & 2.3116 & 0.6566 \\
Beca parcial -- Maestría o superior & 3 & 1.4796 & 2.3116 & 1.5624 \\
\midrule
\multicolumn{4}{r}{$X_0^2 = \sum$} & 4.5678 \\
\bottomrule
\end{tabular}
\end{mitabla}

\textit{Paso 4. Determinar el estadístico teórico.} Con $k = 3$ filas y $m = 2$ columnas, los grados de libertad son $(k-1)(m-1) = (3-1)(2-1) = 2$. De la tabla de chi-cuadrada: $X_{0.05,\,2}^2 = 5.991$.

\textit{Paso 5. Establecer la relación entre estadísticos:} $X_0^2 < X_{0.05,\,2}^2 \;\Rightarrow\; 4.5678 < 5.991$.

\textit{Paso 6. Conclusión.}

\begin{center}
\begin{tabular}{l}
Si ($X_0^2 < X_{0.05,\,2}^2$)\\
\quad \fbox{Independencia}\\
Sino\\
\quad Relación entre las variables\\
FinSi
\end{tabular}
\end{center}

No se puede rechazar que financiamiento y nivel cursado sean independientes.

\textit{Observación:} 4 de 6 celdas tienen $FE < 5$, por lo que la prueba pierde confiabilidad y la conclusión debe tomarse con cautela.

